\documentclass[11pt]{amsart}
\usepackage{amssymb}
\usepackage[margin=1.15in]{geometry}
\usepackage[hidelinks]{hyperref}
\hypersetup{pdftitle={Gabber's theorem in H^2 form}}

\newtheorem{theorem}{Theorem}[section]
\newtheorem{lemma}[theorem]{Lemma}
\newtheorem{proposition}[theorem]{Proposition}
\newtheorem{corollary}[theorem]{Corollary}
\theoremstyle{definition}
\newtheorem{definition}[theorem]{Definition}
\theoremstyle{remark}
\newtheorem{remark}[theorem]{Remark}

\newcommand{\Spec}{\operatorname{Spec}}
\newcommand{\Hom}{\operatorname{Hom}}
\newcommand{\End}{\operatorname{End}}
\newcommand{\Arr}{\operatorname{Arr}}
\newcommand{\Split}{\operatorname{Split}}
\newcommand{\cl}{\operatorname{cl}}
\newcommand{\ev}{\operatorname{ev}}
\newcommand{\Id}{\operatorname{Id}}
\newcommand{\close}{\Diamond}
\newcommand{\tr}[1]{\|#1\|}
\newcommand{\trz}[1]{\|#1\|_0}
\newcommand{\Rx}{R^\times}
\newcommand{\Mat}{M_N(R)}
\newcommand{\code}[1]{\texttt{#1}}
\newcommand{\U}{\mathcal U}
\newcommand{\ap}{\mathsf{ap}}

\title{Gabber's theorem in $H^2$ form, internally}
\date{November 3, 2026}

\newcount\stamphour \newcount\stampmin
\stamphour=\time \divide\stamphour by 60
\stampmin=\time \multiply\stamphour by 60 \advance\stampmin by -\stamphour \divide\stamphour by 60
\edef\stamptime{\the\year-\ifnum\month<10 0\fi\the\month-\ifnum\day<10 0\fi\the\day\space
 \ifnum\stamphour<10 0\fi\the\stamphour:\ifnum\stampmin<10 0\fi\the\stampmin}
\AddToHook{shipout/foreground}{\put(0.5\paperwidth,-\dimexpr\paperheight-0.45in\relax)%
 {\makebox[0pt]{\scriptsize\ttfamily unit/doc/\jobname.tex\quad compiled \stamptime}}}

\begin{document}
\begin{abstract}
In synthetic algebraic geometry with the root topology (a refinement of the
\'etale topology, Section~1), we prove: for a finitely presented $A$ and a
family $G$ of small $\Rx$-gerbes over $\Spec A$ which are stacks, if
$k\cdot\cl(G)=0$ in the group $H^2$ for some $k\geq1$, then there merely is
an Azumaya algebra $E$ over $\Spec A$ with $\delta(E)=\cl(G)$ in $H^2$. The gerbe form, an Azumaya algebra $E$ whose
splittings are fully faithfully compared with $G$, was obtained before. Here
we add the splitting gerbe $\Split(E)$, show that it is a stack and a gerbe,
define $H^2$ as a set truncation of maps into stack gerbes, and show that
the comparison $G\to\Split(E)$ is an equivalence of banded stacks, so that
the classes are equal by univalence. The group law on $H^2$ is the
contracted product, built from families of rank-one modules of a given
weight; $k\cdot\cl(G)=0$ holds exactly when the band change $k_*G$ has a
global object. The sheafification $\close$ is
constructed (by towers of root adjunctions), not assumed, and the splitting
gerbe is made small by replacement, so that the gerbes of the hypothesis and
those of $H^2$ live in one universe. The whole argument is checked in Unit.
\end{abstract}
\maketitle

\section{Setting}

We work in univalent type theory with function extensionality, propositional
truncation $\tr{-}$ and set quotients (both obtained from replacement), over
a local ring $R$ satisfying the axioms of synthetic algebraic geometry.
Universes are $\U_0\subset\U_1$; $R$ and modules live in
$\U_0$.

\emph{Topology.} For a monic polynomial $f$ of positive degree let
$G_f=\tr{\{x\mid f(x)=0\}}$, a proposition. (Degree $0$ is excluded: $f=1$
would give $G_f$ empty and make every sheaf contractible.) We call the
topology generated by the $G_f$ the \emph{root topology}. It is the one used
in~\cite{sevb}, Remark~1.1.2: the \'etale topology is generated by the
$G_f$ with $f$ unramifiable in the sense of Wraith, and this condition is
never used, so all monic $f$ of positive degree are allowed. The root
topology is therefore finer than the \'etale topology; it contains, e.g.,
the root of $T^p-a$ in characteristic $p$, which is not an \'etale cover.
Everything below, including the definition of $H^2$, is stated and checked
for the root topology.

A type $X$ is a \emph{sheaf} if for every $f$ the constant
map $X\to(G_f\to X)$ is an equivalence. Sheaves are closed under $\Pi$,
$\Sigma$, identity types and retracts, and $R$ is a sheaf. For a type $P$ we
write $\close P$ for the sheafification of $\tr P$ (``$P$ holds
locally''), constructed as follows.

A \emph{tower} over a type $Y$ is either empty, with points $Y$, or a
monic polynomial $c_y$ of positive degree depending on $y:Y$ followed by a
tower over $\sum_{y:Y}\{x\mid c_y(x)=0\}$; its points are those of the last
stage, each lying over a point of $Y$. For $P\in\U_0$ let
\[\close P=\tr{\text{a tower over }1\text{ all of whose points give }\tr P}.\]

\begin{proposition}[the closure]\label{p:close}
$\close P$ is a proposition and a sheaf, $P\to\close P$, and $\close P\to Q$
for every proposition $Q$ (in any universe) which is a sheaf and receives
$P$. The sheaves in $\U_0$ form a sheaf.
\end{proposition}
\begin{proof}
$P\to\close P$ by the empty tower. \emph{Elimination}: by induction on the
tower, a map from its points to $Q$ gives a map from its base to $Q$; at a
stage $c$, for each $y$ a map from the roots of $c_y$ to $Q$ factors through
$G_{c_y}$, and $Q$ is a sheaf. \emph{Sheaf}: given $G_f\to\close P$, the
points of $\Spec B$, $B=R[X]/(f)$, are the roots of $f$, so at each point of
$\Spec B$ there is merely a tower; by $Z$-choice and the flattening of
towers (locally constant shapes, clearing of denominators, products) there
is merely one tower over $\Spec B$ with the same property, and prefixed by
the stage $f$ it is a tower over $1$. \emph{Sheaves form a sheaf}: a family
$X:G_f\to\mathrm{Sheaf}$ glues to $\prod_{p:G_f}X_p$, a sheaf, and since
$G_f$ is a proposition evaluation at $p$ is an equivalence
$\prod_{p}X_p\simeq X_p$, so by univalence the glued sheaf restricts to $X$;
a constant family $X$ glues to $G_f\to X$, equal to $X$ since $X$ is a
sheaf.
\end{proof}

Two facts are used repeatedly:
\begin{enumerate}
\item a sheaf proposition that holds locally holds;
\item if $P$ is a proposition and $G_f\to P$ with $P$ a sheaf, then $P$.
\end{enumerate}

\emph{Modules.} A \emph{basis} of rank $N$ of a module $M$ is a family
$(e_i,c_i)_{i<N}$ of vectors and linear forms with $c_i(e_j)=\delta_{ij}$
and $\sum_i c_i(x)e_i=x$. Equivalently, a family $w$ such that
$a\mapsto\sum a_iw_i$, $R^N\to M$, is an equivalence. We use

\begin{lemma}\label{l:free-sheaf}
If $M$ merely has a basis, the carrier of $M$ is a sheaf.
\end{lemma}
\begin{proof}
With a basis, $M$ is a retract of $R^N$ ($x\mapsto(c_i(x))$,
$a\mapsto\sum a_ie_i$), and $R^N$ is a sheaf. For a set, being a sheaf
holds as soon as it merely holds.
\end{proof}

\begin{theorem}[descent of freeness]\label{t:D}
A module whose carrier is a sheaf and which locally has a basis of rank $N$
merely has a basis of rank $N$.
\end{theorem}

This is Theorem~D of~\cite{sevb} (1.3.5), for the root topology,
proved earlier; we use it as a black box.

\section{Gerbes, stacks and $H^2$}

\begin{definition}
A \emph{band} at a point $x$ of a $1$-type $X$ is a group isomorphism
$\beta_x:(x=x)\to\Rx$ (a map with contractible fibres turning composition
of loops into products). A \emph{stack gerbe} is a type $X\in\U_0$ with a
band at every point, such that
\begin{itemize}
\item $X$ is a $1$-type, $\close X$, and $\close(a=b)$ for all $a,b$;
\item each identity type $a=b$ is a sheaf;
\item $X$ itself is a sheaf (\emph{descent for objects}).
\end{itemize}
A \emph{gerbe} is the same without the last condition; a gerbe whose
carrier is a sheaf is a stack gerbe.
\end{definition}

Because the paths of a gerbe are sheaves and locally inhabited, they form
an $\Rx$-torsor which is a sheaf, hence (Hilbert~90) any two objects are
merely equal.

\begin{proposition}[identity of stack gerbes]\label{p:sip}
Let $G,G'$ be stack gerbes and $e:G\simeq G'$ an equivalence of carriers
preserving the bands: $\beta'_{e(x)}(\ap_e(\ell))=\beta_x(\ell)$. Then
$G=G'$.
\end{proposition}
\begin{proof}
By induction on equivalences (univalence in $\U_0$) we may take $e$ to be
the identity. Then the bands agree pointwise ($\ap_{\mathrm{id}}\ell=\ell$),
and the remaining data are propositions: being a $1$-type, $\close X$,
$\close(a=b)$, being a sheaf, and, for a given map
$(x=x)\to\Rx$, having contractible fibres and being multiplicative.
\end{proof}

\begin{definition}
For a type $X$ (here $X=\Spec A$), $H^2(X)=\trz{X\to\mathrm{StackGerbe}}\in\U_1$,
the quotient of $X\to\mathrm{StackGerbe}$ by mere equality. For
$G:X\to\mathrm{StackGerbe}$ we write $\cl(G)$ for its class.
\end{definition}

By Proposition~\ref{p:sip} and function extensionality, $\cl(G)=\cl(G')$ as
soon as there is, at each $x$, a band-preserving equivalence $G_x\simeq
G'_x$. By Theorem~\ref{t:k2}, $\mathrm{StackGerbe}$ is the
sheafification of $K(\Rx,2)$, so $H^2(X)=\trz{X\to\bigcirc K(\Rx,2)}$ is
the second cohomology in the sense of homotopy type theory.

\section{The splitting gerbe}

Fix $N=n+1\geq1$ and an $R$-algebra $E$.

\begin{definition}
A \emph{splitting} of $E$ is a module $M$ which merely has a basis of rank
$N$, with an algebra isomorphism $\sigma:E\cong\End M$. An \emph{arrow}
$(M,\sigma)\to(M',\sigma')$ is a linear bijection $a:M\to M'$ with
$a(\sigma(s)x)=\sigma'(s)(a x)$ for all $s\in E$. Arrows are determined by
their underlying map. $\Split(E)\in\U_1$ is the type of splittings (its objects are modules).
\end{definition}

\begin{proposition}\label{p:paths}
The map $(o=o')\to\Arr(o,o')$ sending $\mathsf{refl}$ to the identity is an
equivalence. Hence $\Split(E)$ is a $1$-type.
\end{proposition}
\begin{proof}
It suffices that $\sum_{o'}\Arr(o,o')$ is contractible. An element is a
module $M'$ with a linear equivalence $e:M\simeq M'$, a mere basis and a
splitting $\sigma'$ of $M'$ intertwined by $e$. By the structure identity
principle for modules we may take $e=\mathrm{id}$. Then
$\sigma(s)x=\sigma'(s)x$ for all $s,x$, so $\sigma=\sigma'$; the bases
agree (a proposition), and the arrow is the identity.
\end{proof}

\begin{proposition}[essentially small]\label{p:small}
Let $D=\{\sigma: E\cong\End(R^N)\}\in\U_0$ and $f:D\to\Split(E)$,
$f(\sigma)=(R^N,\sigma)$. The image $\Split_0(E)\in\U_0$ of $f$ given by
replacement comes with an embedding $j:\Split_0(E)\to\Split(E)$, and $j$ is
an equivalence.
\end{proposition}
\begin{proof}
Replacement applies since $\Split(E)$ is locally small: $o=o'$ is
equivalent to $\Arr(o,o')\in\U_0$ (Proposition~\ref{p:paths}). The fibres
of the embedding $j$ are propositions; they are inhabited: given
$o=(M,\sigma)$ we may choose a basis of $M$, i.e.\ $q:M=R^N$, and
$(R^N,q_*\sigma)=(M,\sigma)$ by induction on $q$; so $o$ is $f(q_*\sigma)$.
\end{proof}

\begin{lemma}[centre]\label{l:centre}
Let $M$ have a basis of rank $N\geq1$ and let $\gamma$ be a linear bijection
of $M$ commuting with every endomorphism. Then $\gamma=\mu$ for a unique
unit $\mu$.
\end{lemma}
\begin{proof}
Let $f_i=c_0(-)e_i$. Then $\gamma(e_i)=\gamma(f_ie_0)=f_i(\gamma
e_0)=\mu e_i$ with $\mu=c_0(\gamma e_0)$, so $\gamma=\mu$ by linearity and
$\sum c_i(x)e_i=x$. If $\gamma y=e_0$ then $1=c_0(\mu y)=\mu c_0(y)$.
Uniqueness: compare at $e_0$.
\end{proof}

\begin{proposition}[band]\label{p:band}
For a splitting $o$, the scalars give an equivalence
$\Rx\simeq\Arr(o,o)$, and the composite
$\beta_o:(o=o)\simeq\Arr(o,o)\simeq\Rx$ is a band.
\end{proposition}
\begin{proof}
An arrow $a$ of $o$ commutes with every endomorphism $f$ of $M$, since
$f=\sigma(s)$ for $s=\sigma^{-1}(f)$; apply Lemma~\ref{l:centre} (the goal
is a proposition, so we may use a basis). Multiplicativity: the arrow of a
composite path is the composite of the arrows (by path induction), and
the composite of the scalars $u$ and $v$ is $uv$.
\end{proof}

\begin{theorem}\label{t:split}
If the carrier of $E$ is a sheaf and locally $E\cong\Mat$, then
$\Split_0(E)$ is a stack gerbe.
\end{theorem}
\begin{proof}
\emph{Paths are sheaves.} By Lemma~\ref{l:free-sheaf} the carriers of $M,M'$
are sheaves, so $\Arr(o,o')$ is a sheaf: linear maps $M\to M'$ form a sheaf
($\Pi$ and identity types), being a bijection is a retract of the data
$(g,h,fg=1,hf=1)$, a sheaf, and intertwining is a $\Pi$ of identities. By
Proposition~\ref{p:paths}, $o=o'$ is a sheaf.

\emph{Locally inhabited.} $R^N$ with its standard basis gives
$\End(R^N)\cong\Mat$; where $E\cong\Mat$ this is a splitting of $E$, an
element of $D$.

\emph{Locally connected.} Two splittings $(M,\sigma)$, $(M',\sigma')$ are
merely isomorphic: choose bases, so that both give automorphisms of $\Mat$;
by Skolem--Noether over the local ring $R$, $\sigma'\sigma^{-1}$ is merely
a conjugation, which yields an arrow.

\emph{Band.} Proposition~\ref{p:band}.

\emph{Descent for objects.} Let $f:G_g\to\Split(E)$ ($G_g$ a generator, a
proposition). Put
\[M=\prod_{p:G_g}M_{f(p)},\qquad \sigma(s)(m)(p)=\sigma_{f(p)}(s)(m(p)).\]
This is the key construction: since $G_g$ is a proposition, for each $p$
evaluation at $p$ is a linear bijection $M\cong M_{f(p)}$ carrying $\sigma$ to $\sigma_{f(p)}$. Hence:
$M$ has a basis under the hypothesis $p:G_g$; its carrier is a sheaf (a
product of sheaves); so by Theorem~\ref{t:D} it merely has a basis.
Likewise $\sigma$ is an equivalence under $p$, and being an equivalence
between sheaves is a sheaf proposition, so $\sigma$ is an equivalence. The
glued splitting $o_f=(M,\sigma)$ restricts to $f$: evaluation at $p$ is an
arrow $o_f\to f(p)$, hence a path (Proposition~\ref{p:paths}). Conversely,
for a constant family $f=\lambda p.\,o$, the constant map
$M_o\to\prod_{p}M_o$ is a bijection because $M_o$ is a sheaf, and it is an
arrow $o\to o_f$. So $o\mapsto\lambda p.\,o$ is an equivalence.

\emph{The small model.} All of this passes to $\Split_0(E)$ along the
equivalence $j$ (Proposition~\ref{p:small}): its identity types are
retracts of the arrows, it is a sheaf, and $\beta_a(\ell)=\beta_{j
a}(\ap_j\ell)$ is a band since $\ap_j$ is an equivalence preserving
composition.
\end{proof}

\section{Gabber's theorem}\label{s:gabber}

\emph{Twisted modules.} For a gerbe $G$ and a homomorphism
$\chi:\Rx\to\Rx$, a family $V:G\to\mathrm{Mod}$ has \emph{weight $\chi$} if
each loop $\ell$ at $g$ acts on $V_g$ (by transport) as multiplication by
$\chi(\beta_g\ell)$. The band change $\chi_*G$ is the type of weight-$\chi$
families of modules whose fibres merely have a basis of rank one; it is a
stack gerbe banded by $\Rx$ (Proposition~\ref{p:lines}). Here we use its
objects: a weight-$k$ family of rank-one modules on $G$. By
Section~\ref{s:group}, $k_*G_x$ has an object for all $x$, merely and
uniformly in $x$, exactly when $k\cdot\cl(G)=0$.

\begin{theorem}[gerbe form, proved before]\label{t:gerbe}
Let $A$ be finitely presented and $G$ a family of gerbes over $\Spec A$
with an element of $\prod_{x:\Spec A}k_*G_x$. Then merely there are $N\geq2$
and a weight-$1$ family $V_x$ on $G_x$ whose fibres merely have bases of
rank $N$, such that the descended endomorphism algebra
$E_x=\End(V_x)^{G_x}$ (invariant sections) is Azumaya (a sheaf algebra,
locally isomorphic to $\Mat$), evaluation at $g$ is an algebra isomorphism
$\ev_g:E_x\cong\End(V_{x,g})$, and for all $g,g'$ the map
\[\varphi:(g=g')\to\{\text{bijections } V_{x,g}\to V_{x,g'}\text{
intertwining }\ev_g,\ev_{g'}\}\]
(transport) is an equivalence, a loop going to its band scalar.
\end{theorem}

This is the result of the flattening, push-forward and normalisation steps
(Gabber~\cite{Gabber}, Ch.~II); we use it as given here, and compare its
proof with the literature in Section~\ref{s:lit}. Its last clause says that
$\Phi_x:G_x\to\Split(E_x)$, $g\mapsto(V_{x,g},\ev_g)$, is fully faithful. We
now upgrade it to an equivalence.

\begin{proposition}\label{p:phi}
Let $G$ be a banded groupoid, $V$ and $E$ as in Theorem~\ref{t:gerbe}, and
$\Phi:G\to\Split(E)$, $\Phi(g)=(V_g,\ev_g)$.
\begin{enumerate}
\item The arrow of $\ap_\Phi(p)$ is $\varphi(p)$; hence $\Phi$ is an
embedding.
\item $\Phi$ preserves the bands.
\item If $G$ is a sheaf and $\close G$, then $\Phi$ is an equivalence, and
so is $j^{-1}\Phi:G\simeq\Split_0(E)$, which preserves the bands.
\end{enumerate}
\end{proposition}
\begin{proof}
(1) By path induction, both are the identity at $\mathsf{refl}$. So
$\ap_\Phi$ is $\varphi$ followed by the inverse of the equivalence of
Proposition~\ref{p:paths}, an equivalence.
(2) $\beta_{\Phi g}(\ap_\Phi\ell)$ is the scalar of $\varphi(\ell)$, which
is $\beta_g(\ell)$.
(3) The fibre $\sum_g(\Phi g=o)\in\U_1$ is a proposition by (1), and a
sheaf ($\Sigma$ of the sheaf $G$ and the sheaves $\Phi g=o$,
Theorem~\ref{t:split}). For $g:G$ it holds, since $\Phi g$ and $o$ are
merely isomorphic (Skolem--Noether, as in Theorem~\ref{t:split}); as
$\close G$, it holds (Proposition~\ref{p:close}, eliminating into a sheaf
proposition of $\U_1$). For the bands of $j^{-1}\Phi$: $\ap_j\ap_{j^{-1}\Phi}=
\ap_{jj^{-1}\Phi}$, and $jj^{-1}\Phi=\Phi$; by induction on this path the
bands are those of $\Phi$, which are those of $G$ by (2).
\end{proof}

\begin{theorem}[$H^2$ form, object hypothesis]\label{t:main}
Let $A$ be finitely presented and $G:\Spec A\to\text{gerbes}$ such
that each $G_x$ is a sheaf (so $G$ is a family of stack gerbes), with an
element of $\prod_{x:\Spec A}k_*G_x$ for some $k\geq1$. Then merely there are $N$ and
Azumaya algebras $E:\Spec A\to\mathrm{AZ}_N$ with
\[\delta(E)=\cl(G)\quad\text{in } H^2(\Spec A),\]
where $\delta(E)=\cl(x\mapsto\Split_0(E_x))$.
\end{theorem}
\begin{proof}
The goal is a proposition, so we may use the data of
Theorem~\ref{t:gerbe}. Each $\Split_0(E_x)$ is a stack gerbe
(Theorem~\ref{t:split}). By Proposition~\ref{p:phi}, $j^{-1}\Phi_x:G_x\simeq
\Split_0(E_x)$ is a band-preserving equivalence, so $G_x=\Split_0(E_x)$ by
Proposition~\ref{p:sip}, and the classes agree.
\end{proof}

\section{The group $H^2$ and torsion classes}\label{s:group}

\emph{Line families.} A \emph{charged type} is a type $X\in\U_0$ with
$\close X$, any two points merely equal, and a function $w$ from loops of
$X$ to $\Rx$ (not required to be a band). $L(X,w)$ is the type of families
$V:X\to\mathrm{Mod}$ whose fibres merely have a basis of rank one and on
which each loop $\ell$ acts by multiplication by $w(\ell)$. A stack gerbe
$G$ is a charged type with $w=\beta$; the band change of
Section~\ref{s:gabber} is $k_*G=L(G,\beta^k)$.

\begin{proposition}\label{p:lines}
If $L(X,w)$ has an object over each point of $X$, then it is (equivalent to)
a stack gerbe.
\end{proposition}
\begin{proof}
Paths of $L(X,w)$ are families of linear bijections (the total space is
contractible), so $L(X,w)$ is a $1$-type. An automorphism of $V$ is a
scalar on each fibre; the scalars are equal along paths, and $X$ is
connected, so they glue (in the sheaf $R$, using $\close X$) to one unit:
this is the band. The isomorphisms $V\cong V'$ form an $\Rx$-torsor which is
a sheaf; it is locally inhabited, since at a point $x$ a bijection
$V_x\cong V'_x$ extends to an invariant section of $\Hom(V,V')$ (weight
$1$, so it descends), hence (Hilbert~90) any two objects are merely equal.
Descent for objects: a family $f$ over a generator $G_i$ (a proposition)
glues to $x\mapsto\prod_{g:G_i}f(g)_x$, which has weight $w$ (transport is
pointwise) and fibres of rank one locally, hence merely
(Theorem~\ref{t:D}). Size: rank-one modules form a locally small type
whose objects are merely $R$, so by replacement it has a small model
$\mathrm{Line}_0$, and families $X\to\mathrm{Line}_0$ with the weight
condition give a small type mapping onto $L(X,w)$; by replacement again
$L(X,w)$ has a small model, and the structure transfers along it.
\end{proof}

\begin{proposition}[maps]\label{p:maps}
A band-preserving map $e:G\to G'$ of stack gerbes is an equivalence, so
$G=G'$.
\end{proposition}
\begin{proof}
On loops $\beta'\circ\ap_e=\beta$ with $\beta,\beta'$ equivalences. For
$a,b$, being an equivalence $\ap_e:(a=b)\to(ea=eb)$ is a sheaf proposition,
and $\close(a=b)$, so by path induction it reduces to loops. The fibres of
$e$ are then propositions and sheaves, and locally inhabited, hence
inhabited. Conclude by Proposition~\ref{p:sip}.
\end{proof}

Three maps produce all comparisons below.
(1) \emph{Pull-back}: if $\varphi:X\to Y$ satisfies
$w_Y(\ap_\varphi\ell)=w_X(\ell)$, then $V\mapsto V\circ\varphi$ maps
$L(Y,w_Y)\to L(X,w_X)$ and keeps the bands (transport along $\ell$ in
$V\circ\varphi$ is transport along $\ap_\varphi\ell$ in $V$).
(2) \emph{Evaluation}: $x\mapsto(V\mapsto V_x)$ maps $X$ to
$L(L(X,w),\beta)$ and carries $w$ to the band.
(3) \emph{Points}: if $k:K$, then $x\mapsto\mathrm{Line}(k=x)$ maps $K$ to
the trivial gerbe $T=L(1,1)$ (rank-one modules) and keeps the bands.

\begin{definition}
$DG=L(G,\beta)$; $G+H=D\,L(G\times H,\beta\times\beta)$, where
$(\beta\times\beta)(\ell)=\beta_G(\ell_1)\beta_H(\ell_2)$; $0\cdot G=T$,
$(n+1)\cdot G=G+n\cdot G$.
\end{definition}

(The twisted line $\mathrm{Line}(s=-)$ is an object of $DG$, and
$s\mapsto\mathrm{Line}(s=-)$ inverts the band, so $D$ reverses the class;
this is why the dual appears in the sum.)

\begin{theorem}[group law]\label{t:group}
For stack gerbes, $DDG=G$, $G+H=H+G$, $(G+H)+K=G+(H+K)$, $G+T=G$,
$G+DG=T$, and $n\cdot G=D(n_*G)$. Hence $H^2(X)$, with
$\cl(G)+\cl(H)=\cl(x\mapsto G_x+H_x)$, $-\cl(G)=\cl(DG)$ and
$0=\cl(x\mapsto T)$, is an abelian group, and
\[n\cdot\cl(G)=\cl(x\mapsto D(n_*G_x)).\]
\end{theorem}
\begin{proof}
$DDG=G$ by evaluation, and a gerbe with a point is $T$ by (3); $DT$ has the
point $V\mapsto V_\ast$, so $DT=T$. Commutativity: pull back along the swap.
Unit: $g\mapsto(g,\ast)$ carries $\ell$ to $(\ell,\mathsf{refl})$, so
$L(G\times T)=DG$ and $G+T=DDG=G$. Inverse: $(g,V)\mapsto V_g$ is a point of
$L(G\times DG,\beta\times\beta)$, so it is $T$ and $G+DG=DT=T$.
Associativity: both $L((G+H)\times K)$ and $L(G\times(H+K))$ are pull-backs
to $G\times H\times K$ (with weight $\beta\beta\beta$) along a product of an
evaluation and an identity. Multiples: the constant family $R$ is a point of
$0_*G$; and $g\mapsto(g,\ev g)$, $G\to G\times D(n_*G)$, carries $\ell$ to
$\beta(\ell)\beta(\ell)^n$, so $L(G\times D(n_*G))=(n+1)_*G$ and
$(n+1)\cdot G=G+D(n_*G)=D((n+1)_*G)$. The operations respect mere equality,
so they descend to the set quotient, and the laws, being propositions, are
checked on classes.
\end{proof}

\begin{proposition}\label{p:zero}
$\cl(K)=0$ if and only if $\tr{\prod_{x:X}K_x}$.
\end{proposition}
\begin{proof}
The quotient is effective, so $\cl(K)=0$ gives $\tr{K=(x\mapsto T)}$, a
mere path of maps, along which the point of $T$ moves to a global section.
Conversely a global section makes each $K_x$ trivial by (3).
\end{proof}

The truncation surrounds the whole family: $\prod_x\tr{K_x}$ would not
suffice.

\begin{theorem}[Gabber, $H^2$ form]\label{t:torsion}
Let $A$ be finitely presented and $G:\Spec A\to\text{gerbes}$ with each
$G_x$ a sheaf. If $k\cdot\cl(G)=0$ in $H^2(\Spec A)$ for some $k\geq1$,
then merely there are $N$ and $E:\Spec A\to\mathrm{AZ}_N$ with
$\delta(E)=\cl(G)$.
\end{theorem}
\begin{proof}
By Theorem~\ref{t:group}, $\cl(x\mapsto D(k_*G_x))=0$, so by
Proposition~\ref{p:zero} merely each $D(k_*G_x)$ has a point, uniformly in
$x$. A point makes $D(k_*G_x)$ trivial, so $k_*G_x=DD(k_*G_x)=DT=T$ has a
point. This is the hypothesis of Theorem~\ref{t:main}, and the goal is a
proposition.
\end{proof}

\subsection{Stack gerbes are the sheafification of \texorpdfstring{$K(\Rx,2)$}{K(Rx,2)}}\label{s:k2}

Let $K$ be the type of $\Rx$-gerbes in the plain sense: $1$-types $X\in\U_0$
with $\tr X$, $\tr{a=b}$ for all $a,b$, and a band at every point. Let $T_0$
be the gerbe of $\Rx$-torsors (made small by replacement), with point
$k_0$ the trivial torsor.

\begin{proposition}\label{p:k2}
$K$ is $K(\Rx,2)$: any two points of $K$ are merely equal, and
$(T_0=T_0)\simeq\mathrm{Tor}(\Rx)$, the type of $\Rx$-torsors.
\end{proposition}
\begin{proof}
As in Proposition~\ref{p:sip}, a band-preserving equivalence $e:G\simeq G'$
gives $G=G'$, and if moreover $e(g)=g'$ then $(G,g)=(G',g')$ as pointed
gerbes. If $g:G$, then $y\mapsto(g=y)$ is a band-preserving equivalence
$G\simeq T_0$ sending $g$ to $k_0$. Hence any $G$ is merely equal to $T_0$
(the goal is a proposition, so we may take a point), and
$\sum_{G:K}G$ is contractible with centre $(T_0,k_0)$. By the fundamental
theorem of identity types, transport of $k_0$ is an equivalence
$(T_0=G)\simeq G$; at $G=T_0$ this is $(T_0=T_0)\simeq T_0\simeq\mathrm{Tor}(\Rx)$.
\end{proof}

\begin{proposition}\label{p:stack-sheaf}
$\mathrm{StackGerbe}$ is a sheaf.
\end{proposition}
\begin{proof}
It is a retract of $\sum_{Y}D(Y)$, where $Y$ ranges over the sheaves in
$\U_0$ and $D(Y)$ is the band, the $1$-type condition, $\close Y$ and
$\close(a=b)$: the paths of a sheaf are sheaves, and the remaining data are
propositions. The sheaves of $\U_0$ form a sheaf (Proposition~\ref{p:close}).
Each $D(Y)$ is a sheaf, being built by $\Pi$, $\Sigma$ and identity types from
$Y$, its paths, $\Rx$ (a sheaf, as $R$ is) and the sheaf propositions
$\close(-)$. Sheaves are closed under $\Sigma$ and retracts.
\end{proof}

The closure works on small propositions; for a type $F$ in $\U_1$ we say
that $F$ holds \emph{locally} if there is a small $A$ with $\close A$ and a
map $A\to F$. A map $f$ between $2$-types is \emph{$\bigcirc$-connected} if
each fibre $F$ of $f$ holds locally, and so do $x=y$, $p=q$ and
$\alpha=\beta$ for all $x,y:F$, $p,q:x=y$ and $\alpha,\beta:p=q$. This is
$\bigcirc$-connectedness of the fibres for the modality $\bigcirc$ of the
root topology~\cite{RSS}: $\bigcirc$ is topological, hence lex, so it
preserves identity types and $\bigcirc F$ is contractible iff $\close F$ and
each $\bigcirc(x=y)$ is contractible; the paths at depth $3$ are
propositions, for which $\bigcirc$-connected means holding locally.

\begin{theorem}\label{t:k2}
The map $\eta:K\to\mathrm{StackGerbe}$ that weakens $\tr-$ to $\close$ is
defined, and it is $\bigcirc$-connected. Hence $\mathrm{StackGerbe}$ is a
sheaf with a $\bigcirc$-connected map from $K(\Rx,2)$, that is,
$\mathrm{StackGerbe}=\bigcirc K(\Rx,2)$, and
$H^2(X)=\trz{X\to\bigcirc K(\Rx,2)}$, with $\bigcirc$ applied to $K$, not to
the set truncation. Post-composition with $\eta$ induces
$\trz{X\to K}\to H^2(X)$.
\end{theorem}
\begin{proof}
$\eta$ is defined: a gerbe $G$ with a point $g$ maps to the trivial stack
gerbe $T$ of lines by $x\mapsto\mathrm{Line}(g=x)$, keeping the bands (map
(3)); $T$ is a gerbe, so the map is an equivalence, and the carrier of $G$
is a sheaf (being a sheaf is a proposition, so a point may be assumed),
hence also its paths. Next,
$K\simeq\sum_{S:\mathrm{StackGerbe}}\tr{S}$: send $G$ to $(\eta G,\tr G)$,
and $(S,t)$ to the gerbe with carrier $S$, inhabited by $t$, connected by
Hilbert~90. Both round trips change only propositions. So the fibre of
$\eta$ at $S$ is that of the first projection, the proposition $\tr S$. It
holds locally, since $\close S$ and a point of $S$ gives a point of the
fibre; and a proposition that holds locally is $\bigcirc$-connected.
\end{proof}

So $\eta$ is an embedding: the gerbes are exactly the stack gerbes with a
mere point, and $\eta$ adds the stack gerbes that have points only locally.

\section{Relation to the literature}\label{s:lit}

\emph{The strategy is classical.} For affine schemes, $\mathrm{Br}=\mathrm{Br}'$ is
due to Gabber~\cite{Gabber}. Hoobler~\cite{Hoobler80} introduced, for regular
rings, the local-to-global argument showing that the set of $f$ over which the
class is represented is an ideal; Hoobler~\cite{Hoobler} simplified Gabber's
proof, and Lieblich~\cite[\S3.1.4]{Lieblich} simplified it further with twisted
sheaves. De Jong~\cite{deJong}, for schemes with an ample line bundle, takes the
affine case as input (his Lemma~3.4 quotes Hoobler's Theorem~7). The skeleton of
Theorem~\ref{t:gerbe} is Lieblich's: twisted locally free sheaves on the gerbe,
$\End$ at the end; the ideal argument with gluing over $D(f)\cup D(g)$;
reduction of the \'etale-local to the Zariski-local situation by Gabber's
splitting-algebra lemma and push-forward along finite locally free covers
(\cite[3.1.3.5--6]{Lieblich}; our flattening by towers of root adjunctions
plays this role); $K$-theory through the nilpotence of the rank-zero part of
$K_0$ and roots in $K_0\otimes\mathbb Q$. The identification
$\delta(E)=\cl(G)$ of Section~4 is the standard dictionary between gerbes and
splittings; here it is formulated and checked internally.

\emph{The gluing step is different.} The arguments of Gabber,
Knus--Ojanguren and Lieblich use, in the gluing, the following facts about
$K_0(A)$:
\begin{itemize}
\item \emph{effectiveness}: some multiple of a class of positive rank is the
  class of a projective module;
\item \emph{cancellation after taking direct sums}: for finite projective
  modules $P,Q$ of positive rank, $[P]=[Q]$ implies
  $P^{\oplus m}\cong Q^{\oplus m}$ for some $m\geq1$.
\end{itemize}
Equality of classes already gives a stable isomorphism
$P\oplus A^r\cong Q\oplus A^r$ by the definition of $K_0$; this elementary
fact is distinct from the cancellation result used to remove the added
summand after taking multiples.
Specifically:
\begin{itemize}
\item Gabber~\cite{Gabber}, Ch.~II, Lemma~K, proves both by reducing to
  finitely generated $\mathbb Z$-algebras and applying ``Serre's theorem and the
  cancellation theorem'' (Bass); the gluing argument that follows invokes
  Lemma~K(ii) and~(iii).
\item Knus--Ojanguren~\cite{KO}, Lemma~3.2, realise rational $K_0$-classes by
  projective modules and then apply ``le th\'eor\`eme d'effacement de Bass''
  (Bass, Chap.~IX, \S7).
\item Hoobler~\cite{Hoobler80} (regular case) realises classes by modules of
  large rank using Bass, Chap.~IX, 4.1.
\item Lieblich~\cite[3.1.4.3, 3.1.4.4]{Lieblich} uses Bass cancellation and
  Serre splitting, after the same reduction, to realise $n$-th roots in
  $K_0\otimes\mathbb Q$ by bundles before gluing.
\end{itemize}

Our gluing (companion note \code{gabber-recollement.tex}) avoids general
effectiveness and Bass cancellation. It consists of two parts.
\begin{enumerate}
\item \emph{Weak cancellation, proved elementarily.} For twisted bundles of
  positive rank, $X\oplus T\cong Y\oplus T$ implies
  $X^{\oplus2m}\cong Y^{\oplus2m}$ for some $m\geq1$: once $T$ is a
  summand of $X^{\oplus m}$ and of $Y^{\oplus m}$ (by evaluation and
  projectivity), $X^{\oplus(m+1)}\cong X^{\oplus m}\oplus Y$, and
  substitution gives
  $X^{\oplus2m}\cong X^{\oplus m}\oplus Y^{\oplus m}\cong Y^{\oplus2m}$.
  The conclusion is of a known type (it is the form of cancellation in
  Lieblich's Lemma~3.1.4.3(ii)); the difference here is the elementary
  proof, which needs no dimension estimate or stable range.
\item \emph{Goodness as the induction invariant, and an explicit clutching
  correction.} One keeps $\rho(P(V))=1$ in $K_0\otimes\mathbb Q$ for
  $P(V)=V^{\otimes n}\otimes\Lambda^{-1}$, $\rho(F)=[F]/\operatorname{rk}F$. After
  gluing two good bundles, the defect $\tau=\rho(P(V'))-1$ lies in the kernel
  of restriction to both pieces, hence $\tau^2=0$ by a clutching computation
  (Whitehead's identity and Quillen patching, Lombardi--Quitt\'e XVI.3.4). An
  actual bundle $S$ with $\rho(S)=1-\tau/n$ is obtained from an inverse power of
  the transition matrix plus a free summand, and
  $\rho(P(V'\otimes S))=(1+\tau)(1-\tau/n)^n=1$. No arbitrary positive-rank
  rational class ever has to be realised by a bundle; the only class that must
  be realised has a directly constructed representative.
\end{enumerate}
The second part is the substantive simplification: it removes effectiveness,
Bass cancellation, Serre splitting and the reduction to noetherian rings from the
gluing argument. This construction differs from the passages examined in
Gabber, Knus--Ojanguren, Hoobler~(1980) and Lieblich. This comparison
establishes a simplification of those arguments, rather than historical
priority for the construction. The constructive characterisations of
Azumaya algebras in~\cite{CLN} concern a different result.

The result is a constructive, machine-checked internal proof of the gerbe
form stated in Theorem~\ref{t:gerbe}, following the classical strategy.
General effectiveness and Bass cancellation are replaced by elementary weak
cancellation and an explicit clutching correction. The $H^2$ statement of
Theorem~\ref{t:main} follows by the identification of Sections~3--4.

\section{What is formalised in Unit}

Everything above is checked in Unit, in \code{unit/sag/twisted}:
\code{./segc sag/twisted all} gives 418 modules, 0 failed, about 12\,s.
The trees that link the changed files also check: \code{sag/torsion-c}
(418 modules), \code{sag/gabber-end} (383) and \code{sag/split-b} (389). No
new axiom is added; the Theory parameters are univalence, function
extensionality and propositional univalence (both derived from it) and
replacement (images and their surjectivity; they give $\tr{-}$, set
quotients and $\Split_0$), together with the SAG axioms at one universe
level and Zariski locality. The closure $\close$ is constructed
(Proposition~\ref{p:close}); it is not a parameter.

\smallskip
\begin{center}
\begin{tabular}{ll}
Statement & File\\\hline
Lemma~\ref{l:free-sheaf} and bases $\leftrightarrow$ frames & \code{BasisFrames}\\
Theorem~\ref{t:D} & \code{DescentTheoremD} (earlier)\\
Proposition~\ref{p:close} & \code{EtaleClosure}, \code{TowerDescent}, \code{Flattening}\\
Proposition~\ref{p:sip}, $H^2$, $\cl$ & \code{StackGerbes} (\code{Eq}, \code{H2})\\
Proposition~\ref{p:paths} & \code{SplitSIP}\\
Proposition~\ref{p:small} & \code{SplitSmall}\\
Lemma~\ref{l:centre} & \code{EndCentre}\\
Proposition~\ref{p:band} & \code{SplitBand}\\
Theorem~\ref{t:split}, local connectedness & \code{SplitConnected}\\
Theorem~\ref{t:split}, the rest & \code{SplitGerbe}\\
Proposition~\ref{p:lines} & \code{BandChange}, \code{BandChangeGerbe}, \code{LineGerbe}\\
Proposition~\ref{p:maps} and small models & \code{GerbeMaps}\\
maps (1)--(3), $D$, $T$ & \code{GerbeOps}\\
$G+H$ & \code{GerbeSum}\\
Theorem~\ref{t:group}, pointwise laws & \code{GerbeLaws}\\
Theorem~\ref{t:group} on $H^2$, Proposition~\ref{p:zero} & \code{H2Group}\\
Theorem~\ref{t:gerbe} & \code{GabberAzumaya} (earlier)\\
Proposition~\ref{p:phi} & \code{SplitEquivalence}\\
Theorem~\ref{t:main} & \code{GabberH2}\\
Theorem~\ref{t:torsion} & \code{GabberTorsion}\\
Proposition~\ref{p:stack-sheaf} & \code{StackSheaf}\\
Proposition~\ref{p:k2} & \code{GerbeSpace}\\
Theorem~\ref{t:k2} & \code{StackComparison}
\end{tabular}
\end{center}

\smallskip
The final declaration is
\[\code{GabberTorsion.Closed[R,setR,axioms,loc].At[A,fp,G,stk,n0,torsion].theorem}\]
of type $\tr{\sum_n\sum_{E:\Spec A\to\mathrm{AZ}_{n+2}}\delta(E)=\cl(G)}$,
where \code{stk} says that each $G_x$ is a sheaf and \code{torsion} is an
equality $(n_0+1)\cdot\cl(G)=0$ in the group $H^2(\Spec A)$
(\code{H2Group}); it instantiates \code{GabberTorsion.Over}, stated for any
closure at level $0$, with the constructed one. Theorem~\ref{t:main} is
\code{GabberH2.Closed}. Negative tests are kept: composing scalar arrows in
the wrong order, $\Phi$ sending a loop to the inverse of its band,
eliminating $\close P$ into $\tr P$ without a sheaf hypothesis, evaluation
carrying $w$ to the inverse of the band, taking the weight
$w_W/w_V$ of $\Hom(V,W)$ to be $\beta\times\beta$ by computation (it is
so only up to a proof), and taking the second component of
$\ap_{g\mapsto(g,\ast)}\ell$ to be $\mathsf{refl}$ by computation, taking a
loop of $K$ at $T$ to fix $k_0$ by computation, and reading a stack gerbe
as a gerbe ($\close S$ in place of $\tr S$), are all rejected. Theorem~\ref{t:k2}
is \code{StackComparison.Over[R,setR,axioms,loc,sh]}: \code{localisation}
(sheaf and $\bigcirc$-connected, with \code{eta-connected}) and
\code{Classes[X].map}; Proposition~\ref{p:k2} is \code{GerbeSpace.Over}
(\code{connected}, \code{loops}).

\begin{thebibliography}{9}
\bibitem{CTS} J.-L. Colliot-Th\'el\`ene, A. N. Skorobogatov, \emph{The
  Brauer--Grothendieck group}, Springer, 2021.
\bibitem{Gabber} O. Gabber, Some theorems on Azumaya algebras, in \emph{The
  Brauer group}, LNM 844, Springer, 1981, 129--209.
\bibitem{Hoobler} R. T. Hoobler, When is $\mathrm{Br}(X)=\mathrm{Br}'(X)$?, in
  \emph{Brauer groups in ring theory and algebraic geometry}, LNM 917,
  Springer, 1982, 231--244.
\bibitem{sevb} \emph{Ch\^atelet's theorem in synthetic algebraic
  geometry}, draft (\code{sevb.pdf}), October 2026.
\bibitem{Hoobler80} R. T. Hoobler, A cohomological interpretation of Brauer
  groups of rings, \emph{Pacific J. Math.} 86 (1980), 89--92.
\bibitem{CLN} T. Coquand, H. Lombardi, S. Neuwirth, Azumaya algebras and
  Barr's theorem, arXiv:2306.17679.
\bibitem{deJong} A. J. de Jong, A result of Gabber, preprint, 2003.
\bibitem{RSS} E. Rijke, M. Shulman, B. Spitters, Modalities in homotopy type
  theory, \emph{Logical Methods in Computer Science} 16 (2020).
\bibitem{KO} M.-A. Knus, M. Ojanguren, Cohomologie \'etale et groupe de
  Brauer, in \emph{The Brauer group}, LNM 844, Springer, 1981, 210--228.
\bibitem{Lieblich} M. Lieblich, Twisted sheaves and the period-index problem,
  \emph{Compositio Math.} 144 (2008), arXiv:math/0511244.
\end{thebibliography}

\end{document}
